\chapter{Konstruktionen von Halbordnungen}
\label{b12guide:ch:constructions}

\section{Durchschnitte und Vereinigungen}

\begin{theorem}[Nichtleere Durchschnitte]
\label{b12guide:thm:intersections}
Ist $\varnothing\ne\mathcal R\subseteq\PO(P)$, so gilt
\[
 \bigcap\mathcal R\in\PO(P).
\]
Insbesondere ist der Durchschnitt zweier Halbordnungen auf
$P$ wieder eine Halbordnung auf $P$.
Der leere Durchschnitt relativ zu $P\times P$ ist dagegen
$P\times P$; er ist genau dann eine Halbordnung, wenn $P$
höchstens ein Element hat.
\end{theorem}

\begin{theorem}[Aufwärts gerichtete Vereinigungen]
\label{b12guide:thm:directed-union}
Sei $\varnothing\ne\mathcal R\subseteq\PO(P)$.
Für je zwei $R,S\in\mathcal R$ gebe es ein $T\in\mathcal R$
mit $R\cup S\subseteq T$. Dann ist
\[
 \bigcup\mathcal R\in\PO(P).
\]
Die Bedingung gilt insbesondere, wenn von je zwei Mitgliedern
der Familie eines im anderen enthalten ist.
Eine beliebige Vereinigung zweier Halbordnungen muß dagegen
weder antisymmetrisch noch transitiv sein.
\end{theorem}

\section{Produkte}

\begin{definition}[Produktordnung]\label{b12guide:def:product}
Für halbgeordnete Mengen $(P,\leq_P)$ und $(Q,\leq_Q)$ wird
auf $P\times Q$ die \emph{Produktordnung} definiert durch
\[
 (p,q)\leq_{\times}(p',q')
 \quad\Longleftrightarrow\quad p\leq_P p'\ \land\ q\leq_Q q'.
\]
\end{definition}

\begin{theorem}[Halbordnung und Hauptmengen des Produkts]
\label{b12guide:thm:product}
Die Produktordnung ist eine Halbordnung auf $P\times Q$.
Für $(p,q)\in P\times Q$ gilt
\[
\begin{aligned}
 \downarrow_{\times}(p,q)&=(\downarrow_P p)\times(\downarrow_Q q),\\
 \uparrow_{\times}(p,q)&=(\uparrow_P p)\times(\uparrow_Q q).
\end{aligned}
\]
Für $p,p'\in P$ und $q,q'\in Q$ gilt
\[
 [(p,q),(p',q')]_{\times}=[p,p']_P\times[q,q']_Q.
\]
\end{theorem}

\begin{theorem}[Extremalelemente des Produkts]
\label{b12guide:thm:product-extremal}
Für $S\subseteq P$ und $T\subseteq Q$ gilt
\[
\begin{aligned}
 \MinP_{\times}(S\times T)&=\MinP_P(S)\times\MinP_Q(T),\\
 \MaxP_{\times}(S\times T)&=\MaxP_P(S)\times\MaxP_Q(T).
\end{aligned}
\]
Für $(s,t)\in S\times T$ gilt ferner
\[
\begin{aligned}
 (s,t)=\min(S\times T)&\ \Longleftrightarrow\
       s=\min S\ \land\ t=\min T,\\
 (s,t)=\max(S\times T)&\ \Longleftrightarrow\
       s=\max S\ \land\ t=\max T.
\end{aligned}
\]
Die Gleichungen mit $\min$ und $\max$ schließen jeweils die
Existenz der bezeichneten Elemente ein.
\end{theorem}

\begin{theorem}[Lexikographische Halbordnung]
\label{b12guide:thm:lexicographic}
Auf $P\times Q$ definiert
\[
 (p,q)\leq_{\mathrm{lex}}(p',q')
 \quad\Longleftrightarrow\quad
 p<_P p'\ \lor\ (p=p'\ \land\ q\leq_Q q')
\]
eine Halbordnung. Jeder Vergleich der Produktordnung gilt auch
in dieser lexikographischen Ordnung. Weder die Faktoren noch
die entstehende Ordnung müssen total geordnet sein.
\end{theorem}

\section{Disjunkte und geordnete Summen}

\begin{theorem}[Disjunkte Summe]
\label{b12guide:thm:disjoint-sum}
Seien $(P_i)_{i\in I}$ paarweise disjunkte Mengen und
$R_i\in\PO(P_i)$. Dann ist
\[
 P=\bigcup_{i\in I}P_i,\qquad R=\bigcup_{i\in I}R_i
\]
eine Halbordnung auf $P$. Elemente verschiedener Summanden
sind unvergleichbar. Für $x\in P_i$ gilt
\[
 \downarrow_R x=\downarrow_{R_i}x,\qquad
 \uparrow_R x=\uparrow_{R_i}x.
\]
Weiter gilt
\[
 \MinP_R(P)=\bigcup_{i\in I}\MinP_{R_i}(P_i),\qquad
 \MaxP_R(P)=\bigcup_{i\in I}\MaxP_{R_i}(P_i).
\]
Die Aussage schließt die leere Familie ein.
\end{theorem}

\begin{theorem}[Geordnete Summe zweier Halbordnungen]
\label{b12guide:thm:ordered-sum}
Seien $P,Q$ disjunkt, $R\in\PO(P)$ und $S\in\PO(Q)$.
Dann ist
\[
 R\oplus S=R\cup S\cup(P\times Q)
\]
eine Halbordnung auf $P\cup Q$, die auf beiden Summanden die
bisherige Ordnung induziert und jedes Element von $P$ unter
jedes Element von $Q$ stellt. Sind beide Summanden nichtleer,
so gilt
\[
 \MinP_{R\oplus S}(P\cup Q)=\MinP_R(P),\qquad
 \MaxP_{R\oplus S}(P\cup Q)=\MaxP_S(Q).
\]
Ist ein Summand leer, bleibt die Ordnung des anderen übrig.
\end{theorem}

\begin{theorem}[Erweiterung durch unvergleichbare neue Elemente]
\label{b12guide:thm:isolated-extension}
Für $P\subseteq A$ und $R\in\PO(P)$ ist
\[
 R'=R\cup\Delta_{A\setminus P}
\]
eine Halbordnung auf $A$ mit $R'|_P=R$.
Sie ist die kleinste Halbordnung auf $A$, deren Relationsgraph
$R$ enthält. Jedes neue Element ist mit jedem davon verschiedenen
Element von $A$ unvergleichbar.
\end{theorem}
